\chapter{结构体：古希腊掌管分类的神}

\begin{introduction}
  \item 为什么需要结构体：把不同类型的数据打包成一个整体
  \item 结构体的定义、声明、初始化与成员访问
  \item \texttt{typedef} 与嵌套结构体
  \item 结构体数组与结构体作为函数参数
  \item 内存对齐：结构体到底占多少字节
  \item 位域与联合体（拔高）
\end{introduction}

\chaptermotto{「世界不是由同一种东西组成的，程序也不该只用一种类型描述世界。」}

数组很优雅，但它有一个前提：所有元素类型相同。可现实中的一个“学生”包含学号（整数）、姓名（字符串）、成绩（小数），类型各不相同。你当然可以用三个平行数组分别存，但数据一旦多了，维护起来就是灾难。

结构体（structure）就是为此而生：它把若干个不同类型的数据打包成一个整体，给这个整体起一个名字。学会结构体，你就从“处理数据”迈向了“描述对象”。

\section{为什么需要结构体}

假设要管理 3 个学生的信息。用平行数组的写法是这样的：

\begin{codebox}{parallel-arrays.c —— 用平行数组存学生信息}
#include <stdio.h>

int main(void)
{
    int ids[3] = {1001, 1002, 1003};
    char names[3][20] = {"Alice", "Bob", "Carol"};
    double scores[3] = {92.5, 78.0, 85.5};
    int i;

    for (i = 0; i < 3; i++) {
        printf("%d %s %.1f\n", ids[i], names[i], scores[i]);
    }

    return 0;
}
\end{codebox}

这个写法的问题在于：三个数组必须\textbf{一一对应}。如果排序时只交换了 \code{scores}，忘了同步交换 \code{ids} 和 \code{names}，数据就全乱了。新增一个字段（比如年龄）还要再开一个数组，处处要同步。

结构体把“一个学生的所有信息”绑在一起：

\begin{codebox}{struct-student.c —— 用结构体存学生信息}
#include <stdio.h>

struct Student {
    int id;
    char name[20];
    double score;
};

int main(void)
{
    struct Student s;

    s.id = 1001;
    s.score = 92.5;
    /* s.name 是字符数组，赋值方式第 9 章讲，这里先用初始化 */

    printf("%d %.1f\n", s.id, s.score);
    return 0;
}
\end{codebox}

\begin{figure}[htbp]
  \centering
  \begin{tikzpicture}[x=1cm,y=1cm]
    % 平行数组
    \node[font=\small\bfseries] at (0,3.0) {平行数组};
    \node[membox,minimum width=2.4cm] (id0) at (0,2.2) {1001};
    \node[membox,minimum width=2.4cm,right=0mm of id0] (id1) {1002};
    \node[membox,minimum width=2.4cm,right=0mm of id1] (id2) {1003};
    \node[font=\scriptsize,color=black!55,left=1mm of id0] {ids};

    \node[membox,minimum width=2.4cm] (nm0) at (0,1.3) {Alice};
    \node[membox,minimum width=2.4cm,right=0mm of nm0] (nm1) {Bob};
    \node[membox,minimum width=2.4cm,right=0mm of nm1] (nm2) {Carol};
    \node[font=\scriptsize,color=black!55,left=1mm of nm0] {names};

    \node[membox,minimum width=2.4cm] (sc0) at (0,0.4) {92.5};
    \node[membox,minimum width=2.4cm,right=0mm of sc0] (sc1) {78.0};
    \node[membox,minimum width=2.4cm,right=0mm of sc1] (sc2) {85.5};
    \node[font=\scriptsize,color=black!55,left=1mm of sc0] {scores};

    % 结构体数组
    \node[font=\small\bfseries] at (9.2,3.0) {结构体数组};
    \node[draw=structurecolor!70,fill=structurecolor!6,rounded corners=2mm,
          minimum width=3.2cm,minimum height=1.1cm,align=left,font=\scriptsize] (st0) at (9.2,2.2)
      {id: 1001\\name: Alice\\score: 92.5};
    \node[draw=structurecolor!70,fill=structurecolor!6,rounded corners=2mm,
          minimum width=3.2cm,minimum height=1.1cm,align=left,font=\scriptsize,below=2mm of st0] (st1)
      {id: 1002\\name: Bob\\score: 78.0};
    \node[draw=structurecolor!70,fill=structurecolor!6,rounded corners=2mm,
          minimum width=3.2cm,minimum height=1.1cm,align=left,font=\scriptsize,below=2mm of st1] (st2)
      {id: 1003\\name: Carol\\score: 85.5};
    \node[font=\scriptsize,color=black!55,right=2mm of st1] {每人一“包”};
  \end{tikzpicture}
  \caption{平行数组 vs.\ 结构体数组：后者把一个人的信息绑在一起}
  \label{fig:struct-vs-parallel}
\end{figure}

\section{结构体的定义与使用}

\subsection{定义}

\begin{center}
\fbox{\parbox{0.82\textwidth}{
\ttfamily struct 结构体名 \{ \\
\phantom{xxxx}类型 成员1; \\
\phantom{xxxx}类型 成员2; \\
\phantom{xxxx}... \\
\ttfamily \}; \quad /* 注意这个分号 */
}}
\end{center}

\begin{codebox}{struct-def.c —— 定义一个学生结构体}
struct Student {
    int id;             /* 学号 */
    char name[20];      /* 姓名 */
    double score;       /* 成绩 */
};
\end{codebox}

\begin{pitfall}[结构体定义末尾的分号不能少]
\texttt{struct Student \{ ... \};} 最后的分号是定义的一部分，漏掉会导致一连串奇怪的编译错误。这个分号很容易忘，因为函数定义后面不写分号。
\end{pitfall}

\code{struct Student} 是一个\textbf{类型名}，就像 \code{int} 一样。注意：\code{Student} 单独使用不合法，必须写 \code{struct Student}，除非用 \code{typedef} 起别名（见后文）。

\subsection{声明变量与访问成员}

\begin{codebox}{struct-use.c —— 声明、初始化、访问成员}
#include <stdio.h>

struct Student {
    int id;
    char name[20];
    double score;
};

int main(void)
{
    struct Student s1;
    struct Student s2 = {1002, "Bob", 78.0};

    s1.id = 1001;
    s1.score = 92.5;

    printf("s1: %d, %.1f\n", s1.id, s1.score);
    printf("s2: %d, %s, %.1f\n", s2.id, s2.name, s2.score);

    return 0;
}
\end{codebox}

\begin{outbox}{运行结果}
s1: 1001, 92.5
s2: 1002, Bob, 78.0
\end{outbox}

访问成员用点号 \code{.}，读作“的”：\code{s1.id} 就是“s1 的 id”。

\subsection{整体赋值}

和数组不同，\textbf{同类型的结构体变量可以整体赋值}，编译器会把所有成员逐字节复制过去：

\begin{codebox}{struct-assign.c —— 结构体整体赋值}
#include <stdio.h>

struct Point {
    int x;
    int y;
};

int main(void)
{
    struct Point p1 = {3, 4};
    struct Point p2;

    p2 = p1;        /* 合法：把 p1 的所有成员复制给 p2 */

    printf("p2 = (%d, %d)\n", p2.x, p2.y);
    return 0;
}
\end{codebox}

\begin{outbox}{运行结果}
p2 = (3, 4)
\end{outbox}

\subsection{\texttt{typedef}：给类型起别名}

每次写 \code{struct Student} 有点啰嗦。可以用 \code{typedef} 起一个别名：

\begin{codebox}{typedef-struct.c —— typedef 的两种写法}
/* 写法一：先定义，再起别名 */
struct Student {
    int id;
    char name[20];
    double score;
};
typedef struct Student Student;

/* 写法二：定义的同时起别名（更常用） */
typedef struct {
    int x;
    int y;
} Point;
\end{codebox}

之后就可以直接写 \code{Student s;}、\code{Point p;}，和 \code{int a;} 一样自然。

\begin{tipbox}[本书的风格]
本书在需要频繁使用某个结构体时，会用 \code{typedef} 起一个首字母大写的别名；只在定义处出现一次的结构体，可以直接用 \code{struct} 写法。
\end{tipbox}

\subsection{嵌套结构体}

结构体的成员也可以是另一个结构体：

\begin{codebox}{nested-struct.c —— 嵌套结构体}
#include <stdio.h>

typedef struct {
    int year;
    int month;
    int day;
} Date;

typedef struct {
    int id;
    char name[20];
    Date birthday;      /* 成员是另一个结构体 */
    double score;
} Student;

int main(void)
{
    Student s = {1001, "Alice", {2005, 6, 18}, 92.5};

    printf("%s 出生于 %d 年 %d 月 %d 日\n",
           s.name, s.birthday.year, s.birthday.month, s.birthday.day);

    return 0;
}
\end{codebox}

\begin{outbox}{运行结果}
Alice 出生于 2005 年 6 月 18 日
\end{outbox}

访问嵌套成员就一层层点下去：\code{s.birthday.year}。

\section{结构体在内存里长什么样}

结构体的成员按声明顺序依次存放，但\textbf{不一定紧挨着}——编译器会在成员之间插入“填充字节”（padding），让每个成员的地址满足对齐要求。

\begin{figure}[htbp]
  \centering
  \begin{tikzpicture}[x=1cm,y=1cm]
    % 标尺
    \foreach \i in {0,...,15} {
      \node[font=\scriptsize,color=black!45] at (\i*0.9+0.45,1.3) {\i};
    }
    % 成员
    \draw[structurecolor!70,fill=ptrfill] (0,0) rectangle (0.9,-1.2);
    \node[font=\scriptsize] at (0.45,-0.6) {a};
    \draw[black!15,fill=black!6] (0.9,0) rectangle (3.6,-1.2);
    \node[font=\scriptsize,color=black!45] at (2.25,-0.6) {填充};
    \draw[structurecolor!70,fill=structurecolor!12] (3.6,0) rectangle (7.2,-1.2);
    \node[font=\scriptsize] at (5.4,-0.6) {b};
    \draw[structurecolor!70,fill=okfill] (7.2,0) rectangle (14.4,-1.2);
    \node[font=\scriptsize] at (10.8,-0.6) {c};

    \node[font=\scriptsize,color=black!55,below] at (0.45,-1.2) {偏移 0};
    \node[font=\scriptsize,color=black!55,below] at (5.4,-1.2) {偏移 4};
    \node[font=\scriptsize,color=black!55,below] at (10.8,-1.2) {偏移 8};
    \node[font=\scriptsize,color=black!55,below] at (14.4,-1.2) {总大小 16};
  \end{tikzpicture}
  \caption{\texttt{struct \{ char a; int b; double c; \}} 的内存布局（典型 8 字节对齐）}
  \label{fig:struct-padding}
\end{figure}

为什么需要对齐？因为 CPU 读取内存时，按“字”对齐的访问速度最快，某些架构甚至不允许非对齐访问。编译器为了性能，会在成员之间补空。

\begin{codebox}{sizeof-struct.c —— 结构体的大小}
#include <stdio.h>

struct A {
    char a;
    int b;
    double c;
};

struct B {
    char a;
    char b;
    int c;
};

int main(void)
{
    printf("sizeof(struct A) = %d\n", (int)sizeof(struct A));
    printf("sizeof(struct B) = %d\n", (int)sizeof(struct B));
    return 0;
}
\end{codebox}

\begin{outbox}{运行结果（典型 64 位环境）}
sizeof(struct A) = 16
sizeof(struct B) = 8
\end{outbox}

\code{struct A} 的成员一共 $1 + 4 + 8 = 13$ 字节，但因为对齐，实际占 16 字节。\code{struct B} 的两个 \code{char} 挤在一起，后面补 2 字节再放 \code{int}，共 8 字节。

\begin{tipbox}[省内存的小技巧]
把成员按“从大到小”排列，通常能减少填充。例如 \texttt{struct \{double c; int b; char a;\}} 在很多平台上只占 16 字节，而 \texttt{\{char a; int b; double c;\}} 也是 16——但当成员更多时，差别会明显。
\end{tipbox}

\begin{pitfall}[不要把结构体的内存布局当成协议]
虽然你可以用 \code{memcpy} 把结构体直接写进文件或发到网络，但填充字节的内容是不确定的，不同编译器/平台的对齐规则也不同。跨平台传输数据时，应该逐字段序列化，而不是直接发结构体。
\end{pitfall}

\section{结构体数组}

结构体数组就是“每个元素都是一个结构体”的数组：

\begin{codebox}{struct-array.c —— 结构体数组}
#include <stdio.h>

typedef struct {
    int id;
    char name[20];
    double score;
} Student;

int main(void)
{
    Student cls[3] = {
        {1001, "Alice", 92.5},
        {1002, "Bob",   78.0},
        {1003, "Carol", 85.5}
    };
    int i;
    double sum = 0.0;

    for (i = 0; i < 3; i++) {
        printf("%d %-8s %.1f\n", cls[i].id, cls[i].name, cls[i].score);
        sum = sum + cls[i].score;
    }

    printf("平均分 = %.2f\n", sum / 3.0);
    return 0;
}
\end{codebox}

\begin{outbox}{运行结果}
1001 Alice    92.5
1002 Bob      78.0
1003 Carol    85.5
平均分 = 85.33
\end{outbox}

\begin{note}
\texttt{\%-8s} 表示字符串左对齐、占 8 格，用来让姓名列对齐。这个格式符第 9 章会详细讲。
\end{note}

\section{结构体与函数}

结构体可以像普通变量一样传给函数、从函数返回。这里先讲最基本的用法，函数的完整语法见附录~\ref{app:function}。

\subsection{按值传递：函数拿到一份副本}

\begin{codebox}{struct-param.c —— 结构体作为参数}
#include <stdio.h>

typedef struct {
    int id;
    char name[20];
    double score;
} Student;

void print_student(Student s)
{
    printf("%d %s %.1f\n", s.id, s.name, s.score);
}

int main(void)
{
    Student s = {1001, "Alice", 92.5};

    print_student(s);       /* 传递的是一份副本 */
    return 0;
}
\end{codebox}

按值传递时，函数里修改 \code{s} 不会影响 \code{main} 里的原变量。这很安全，但如果结构体很大，复制会浪费时间。

\subsection{按地址传递：函数可以修改原变量}

如果想让函数修改原结构体，或者结构体很大想避免复制，就传地址：

\begin{codebox}{struct-pointer.c —— 传地址修改结构体}
#include <stdio.h>

typedef struct {
    int id;
    double score;
} Student;

void set_score(Student *p, double new_score)
{
    p->score = new_score;       /* 等价于 (*p).score */
}

int main(void)
{
    Student s = {1001, 92.5};

    set_score(&s, 95.0);        /* 传地址 */
    printf("%.1f\n", s.score);
    return 0;
}
\end{codebox}

\begin{outbox}{运行结果}
95.0
\end{outbox}

\code{p->score} 是“指针访问成员”的写法，等价于 \code{(*p).score}。下一章学指针时会详细解释，现在先记住这个符号。

\section{综合例题}

\subsection{例：学生成绩统计}

\begin{example}
有 5 个学生，每个学生有学号、姓名、成绩。输出平均分、最高分学生，以及高于平均分的人数。
\end{example}
\begin{solution}
思路如下：

\begin{codebox}{student-stats.c —— 学生成绩统计}
#include <stdio.h>

typedef struct {
    int id;
    char name[20];
    double score;
} Student;

int main(void)
{
    Student cls[5] = {
        {1001, "Alice", 92.5},
        {1002, "Bob",   78.0},
        {1003, "Carol", 85.5},
        {1004, "Dave",  66.0},
        {1005, "Eve",   95.0}
    };
    int i;
    double sum = 0.0;
    double average;
    int max_index = 0;
    int above = 0;

    for (i = 0; i < 5; i++) {
        sum = sum + cls[i].score;
    }
    average = sum / 5.0;

    for (i = 1; i < 5; i++) {
        if (cls[i].score > cls[max_index].score) {
            max_index = i;
        }
    }

    for (i = 0; i < 5; i++) {
        if (cls[i].score > average) {
            above++;
        }
    }

    printf("平均分 = %.2f\n", average);
    printf("最高分：%s (%.1f)\n",
           cls[max_index].name, cls[max_index].score);
    printf("高于平均分的人数：%d\n", above);

    return 0;
}
\end{codebox}

\begin{outbox}{运行结果}
平均分 = 83.40
最高分：Eve (95.0)
高于平均分的人数：2
\end{outbox}
\end{solution}

\subsection{例：复数运算}

\begin{example}
定义复数结构体，输入两个复数，输出它们的和与积。
\end{example}
\begin{solution}
复数 $a + bi$ 与 $c + di$ 相加为 $(a+c) + (b+d)i$，相乘为 $(ac-bd) + (ad+bc)i$。

\begin{codebox}{complex.c —— 复数运算}
#include <stdio.h>

typedef struct {
    double re;      /* 实部 */
    double im;      /* 虚部 */
} Complex;

int main(void)
{
    Complex x, y, sum, product;

    printf("请输入第一个复数（实部 虚部）：");
    scanf("%lf%lf", &x.re, &x.im);
    printf("请输入第二个复数（实部 虚部）：");
    scanf("%lf%lf", &y.re, &y.im);

    sum.re = x.re + y.re;
    sum.im = x.im + y.im;

    product.re = x.re * y.re - x.im * y.im;
    product.im = x.re * y.im + x.im * y.re;

    printf("和 = %.2f + %.2fi\n", sum.re, sum.im);
    printf("积 = %.2f + %.2fi\n", product.re, product.im);

    return 0;
}
\end{codebox}

\begin{outbox}{运行示例}
请输入第一个复数（实部 虚部）：1 2
请输入第二个复数（实部 虚部）：3 4
和 = 4.00 + 6.00i
积 = -5.00 + 10.00i
\end{outbox}
\end{solution}

\subsection{例：日期与下一天}

\begin{example}
定义日期结构体，输入一个日期，输出它的下一天。要考虑闰年和月份天数。
\end{example}
\begin{solution}
思路如下：

\begin{codebox}{next-day.c —— 下一天}
#include <stdio.h>

typedef struct {
    int year;
    int month;
    int day;
} Date;

int is_leap(int year)
{
    return (year % 4 == 0 && year % 100 != 0) || (year % 400 == 0);
}

int main(void)
{
    Date d;
    int days[13] = {0, 31, 28, 31, 30, 31, 30, 31, 31, 30, 31, 30, 31};

    printf("请输入日期（年 月 日）：");
    scanf("%d%d%d", &d.year, &d.month, &d.day);

    if (is_leap(d.year)) {
        days[2] = 29;
    }

    d.day++;
    if (d.day > days[d.month]) {
        d.day = 1;
        d.month++;
        if (d.month > 12) {
            d.month = 1;
            d.year++;
        }
    }

    printf("下一天：%d-%02d-%02d\n", d.year, d.month, d.day);
    return 0;
}
\end{codebox}

\begin{outbox}{运行示例}
请输入日期（年 月 日）：2024 2 28
下一天：2024-02-29
\end{outbox}

这里用数组 \code{days} 存每月天数，下标 1\textasciitilde12 对应 1\textasciitilde12 月，\code{days[0]} 占位不用。
\end{solution}

\subsection{例：电机参数}

\begin{example}
定义一个电机参数结构体，包含编号、PID 三个参数和输出限幅；输入一组参数并打印。
\end{example}
\begin{solution}
思路如下：

\begin{codebox}{motor.c —— 电机参数结构体}
#include <stdio.h>

typedef struct {
    int id;
    double kp;
    double ki;
    double kd;
    double limit;
} MotorParam;

int main(void)
{
    MotorParam m;

    printf("请输入电机编号与 PID 参数：");
    scanf("%d%lf%lf%lf", &m.id, &m.kp, &m.ki, &m.kd);

    printf("请输入输出限幅：");
    scanf("%lf", &m.limit);

    printf("电机 %d: kp=%.3f, ki=%.3f, kd=%.3f, limit=%.2f\n",
           m.id, m.kp, m.ki, m.kd, m.limit);

    return 0;
}
\end{codebox}

结构体让一组相关参数“打包”在一起，后面写 PID 函数时，传一个 \code{MotorParam} 就够了，不用传五个参数。
\end{solution}

\section{常见错误}

\begin{center}
\begin{tabular}{@{}p{0.38\textwidth}p{0.52\textwidth}@{}}
\toprule
错误 & 问题 \\
\midrule
结构体定义后漏分号 & 后续代码出现莫名其妙的错误 \\
\code{Student s;} 但没用 \code{typedef} & 应写 \code{struct Student s;} \\
\code{s == t} & 结构体不能直接用 \code{==} 比较 \\
\code{s.name = "Bob"} & 数组不能整体赋值，要用 \code{strcpy} \\
\texttt{scanf("\%s", \&s.name)} & \code{s.name} 本身就是地址，不要加 \texttt{\&} \\
忘记给数值成员加 \texttt{\&} & \code{scanf} 写不进去 \\
\bottomrule
\end{tabular}
\end{center}

\begin{pitfall}[结构体比较要逐成员]
\code{if (s1 == s2)} 在 C 里不合法。正确做法是逐成员比较，或者用 \code{memcmp}——但 \code{memcmp} 会把填充字节也比进去，不可靠。最稳妥的是逐成员比较：
\begin{codebox}{struct-compare.c —— 逐成员比较}
int equal(Student a, Student b)
{
    return a.id == b.id &&
           a.score == b.score;   /* 字符串比较需用 strcmp */
}
\end{codebox}
\end{pitfall}

\begin{pitfall}[字符串成员不能直接赋值]
\code{s.name = "Bob";} 是错的，因为 \code{name} 是数组，数组名不能做赋值目标。正确写法是 \code{strcpy(s.name, "Bob");}，第 9 章会讲。
\end{pitfall}

\section{拔高：位域与联合体}

\subsection{位域}

有时结构体成员只需要几个二进制位（例如标志位），用 \code{int} 会浪费空间。C 语言提供了\textbf{位域}：

\begin{codebox}{bitfield.c —— 位域}
#include <stdio.h>

struct Flags {
    unsigned int enable : 1;    /* 占 1 位 */
    unsigned int mode   : 2;    /* 占 2 位，取值 0-3 */
    unsigned int speed  : 5;    /* 占 5 位，取值 0-31 */
};

int main(void)
{
    struct Flags f;

    f.enable = 1;
    f.mode = 3;
    f.speed = 20;

    printf("enable=%u mode=%u speed=%u\n", f.enable, f.mode, f.speed);
    printf("sizeof(struct Flags) = %d\n", (int)sizeof(struct Flags));
    return 0;
}
\end{codebox}

位域在嵌入式寄存器映射中很常见，但它的具体布局（谁在高位、跨字节怎么排）由编译器决定，跨平台代码要谨慎。

\subsection{联合体}

\code{union} 让多个成员\textbf{共用同一块内存}，同一时刻只有一个成员有意义：

\begin{codebox}{union.c —— 联合体}
#include <stdio.h>

union Data {
    int i;
    float f;
    unsigned char bytes[4];
};

int main(void)
{
    union Data d;

    d.i = 0x12345678;
    printf("i = %x\n", d.i);
    printf("bytes = %02x %02x %02x %02x\n",
           d.bytes[0], d.bytes[1], d.bytes[2], d.bytes[3]);
    printf("sizeof(union Data) = %d\n", (int)sizeof(union Data));

    return 0;
}
\end{codebox}

联合体常用于：
\begin{itemize}
  \item 查看一个数据的字节表示（例如上面的小端序演示）；
  \item 节省内存：同一块空间在不同时刻存不同类型；
  \item 协议解析：把一帧数据按不同格式解释。
\end{itemize}

\begin{note}
上面代码里 \code{d.i = 0x12345678} 后读 \code{d.bytes}，属于“用另一个成员读取”，在 C 标准里是允许的（类型双关），但不同电脑的存放顺序可能不同。这里只用于演示，不要依赖它做可移植代码。
\end{note}

\section{本章小结}

\begin{center}
\begin{tabular}{@{}p{0.24\textwidth}p{0.66\textwidth}@{}}
\toprule
主题 & 要点 \\
\midrule
结构体 & 把不同类型的成员打包成一个整体 \\
定义 & \texttt{struct 名字 \{ ... \};}，末尾分号不能少 \\
访问 & 用 \code{.} 访问成员；整体赋值是逐字节复制 \\
\code{typedef} & 给类型起别名，少写 \code{struct} \\
嵌套 & 成员可以是另一个结构体，逐层点访问 \\
数组 & 结构体数组适合管理一批对象 \\
函数 & 按值传递复制一份；传地址用 \code{->} 修改原值 \\
内存 & 成员按声明顺序排列，可能有填充字节 \\
拔高 & 位域省空间；联合体共用内存 \\
\bottomrule
\end{tabular}
\end{center}

\begin{tipbox}[下一章预告]
结构体让我们能“描述对象”，但还有一个问题没解决：如何让函数直接修改调用者的变量？如何让一个变量“指向”另一个变量？如何高效地传递大块数据？答案就是\textbf{指针}——C 语言最锋利、也最容易伤到自己的工具。
\end{tipbox}

\begin{problemset}
  \item 定义一个 \code{Point} 结构体（\code{x}、\code{y}），输入两个点，输出它们的中点。
  \item 定义一个 \code{Rect} 结构体（左上角、右下角），计算面积和周长。
  \item 定义 \code{Student} 结构体，输入 5 个学生信息，输出平均分和最高分学生。
  \item 用结构体数组实现：输入 5 个学生的成绩，按输入顺序输出，并统计及格人数。
  \item 定义 \code{Date} 结构体，判断两个日期哪个更早。
  \item 定义 \code{Complex} 结构体，实现复数的加法、减法、乘法。
  \item 定义一个包含 \code{char}、\code{int}、\code{double} 的结构体，用 \code{sizeof} 观察它的大小，并解释为什么不是各成员大小之和。
  \item 写一个函数，接收 \code{Student} 的地址，把成绩加 5 分（不超过 100）。
  \item 用位域定义一个 8 位的状态字，包含 1 位使能、2 位模式、5 位速度，并打印各字段。
  \item 思考题：为什么结构体可以整体赋值，而数组不可以？
\end{problemset}
